\chapter{Øvelser og korte svar}

\chapterlead{Øvelserne træner den juridiske arbejdsgang. De har ikke ét magisk svar; et godt svar identificerer de relevante spørgsmål og viser, hvilke oplysninger der mangler.}

\section{Øvelse 1: Den forsinkede levering}

En privatperson bestiller en sofa med levering senest 1. juni. Sofaen ankommer 20. juni. Sælgeren siger, at forsinkelsen skyldes en underleverandør. Køberen havde planlagt at bruge sofaen til en fest 5. juni.

\begin{enumerate}
  \item Hvilke faktiske oplysninger skal du afklare?
  \item Hvilken type misligholdelse er der tale om?
  \item Hvilke retsfølger kan være relevante?
  \item Er festens aflysning automatisk et erstatningsberettiget tab?
\end{enumerate}

\textbf{Kort svar:} Der er sandsynligvis tale om forsinkelse, men man skal kontrollere, om datoen var en bindende del af aftalen, om der er forbrugerforhold, og om særlige vilkår gælder. Ophævelse, afslag, erstatning eller krav om levering kan have forskellige betingelser. En skuffelse eller aflyst fest er ikke i sig selv nok til erstatning; tab, ansvar og påregnelighed skal vurderes.

\section{Øvelse 2: Myndighedens manglende høring}

En kommune modtager et notat fra en nabo, som hævder, at en borger driver erhverv fra sin bolig i strid med en tilladelse. Kommunen inddrager tilladelsen uden at sende notatet til borgeren.

\textbf{Kort svar:} Undersøg først, om beslutningen er en afgørelse, om notatet er en væsentlig faktisk oplysning til ugunst, om borgeren kendte det, og om en undtagelse fra partshøring gælder. Undersøg også hjemmel, oplysningsgrundlag, begrundelse og klage. Manglende høring kan have betydning, hvis borgeren kunne have korrigeret eller nuanceret oplysningen.

\section{Øvelse 3: Ruden}

Et barn spiller bold tæt på et hus. Bolden knuser et vindue. Forældrene siger, at barnet ikke kunne forudse, at bolden ville ramme ruden.

\textbf{Kort svar:} Spørg til alder, situation, boldens retning, afstand, tidligere advarsler, tilsyn og almindelige erstatningsregler. Ansvar kræver ikke kun skade, men et relevant ansvarsgrundlag og årsagssammenhæng. Forældres ansvar og eventuel forsikring skal analyseres særskilt.

\section{Øvelse 4: Ansigtsgenkendelse}

Et bibliotek vil bruge ansigtsgenkendelse til at genkende personer, der tidligere har afleveret bøger for sent.

\textbf{Kort svar:} Det rejser spørgsmål om biometriske oplysninger, behandlingsgrundlag, formål, nødvendighed, proportionalitet, dataminimering, sikkerhed, opbevaring, registreredes rettigheder og eventuelle forbud eller særlige AI-regler. Et praktisk formål gør ikke automatisk indgrebet lovligt.

\section{Øvelse 5: Den ufuldstændige kontrakt}

En freelancer skriver ``jeg kan levere hjemmesiden for cirka 20.000 kr.'' Kunden svarer ``fint, lad os gøre det''. Efter arbejdet sender freelanceren en faktura på 45.000 kr.

\textbf{Kort svar:} ``Cirka'' og arbejdsbeskrivelsen skal fortolkes i aftalens kontekst. Var der et fast tilbud, et overslag eller blot en invitation til at forhandle? Hvilke funktioner var omfattet, og blev ændringer aftalt? Kundens svar kan vise accept, men måske ikke af alle detaljer. Dokumentation af beskeder og projektets omfang er central.

\section{Øvelse 6: Strafferetlig skyld}

En person tager en jakke fra en stol i en café, fordi personen tror, at den er dennes egen. Jakken tilhører en anden.

\textbf{Kort svar:} Den objektive handling og ejerskiftet er ikke alene nok. Undersøg, hvad personen vidste og ville på tidspunktet. En faktisk fejl kan påvirke forsættet, men den præcise betydning afhænger af den relevante bestemmelse og eventuel uagtsomhed.

\section{Øvelse 7: Frist}

En borger modtager en afgørelse digitalt fredag kl. 16.00. Klagefristen er fire uger. Borgeren sender klagen fire uger senere fredag kl. 16.05.

\textbf{Kort svar:} Beregn fristen efter den konkrete klageregel og reglerne om digital meddelelse, helligdage, fristens udløb og modtagelse. Man må ikke gætte ud fra almindelig kalenderlogik. Hvis vejledningen er uklar eller myndighedens system har problemer, kan det få betydning, men det skal dokumenteres.

\section{Øvelse 8: EU-forbindelsen}

En dansk myndighed henviser til EU-charteret i en sag, der kun handler om en rent intern kommunal procedure uden forbindelse til EU-regulering.

\textbf{Kort svar:} Først skal man undersøge, om myndigheden gennemfører eller anvender EU-ret. Charteret gælder ikke som en generel dansk rettighedskatalog for alle situationer. Grundloven og ECHR kan stadig være relevante gennem deres egne anvendelsesregler.

\section{Øvelse 9: Beviset}

En virksomhed siger, at en kunde mundtligt accepterede en ekstra betaling. Kunden benægter det. Der findes kun en intern note skrevet af virksomhedens sælger.

\textbf{Kort svar:} Den interne note kan være bevis, men dens vægt afhænger af tidspunkt, detaljer, sælgerens interesse og øvrige omstændigheder. Virksomheden skal identificere aftalens indhold og bevisbyrden. En intern notits er ikke automatisk bevis for, at kunden accepterede.

\section{Øvelse 10: Det bedste modargument}

Skriv et juridisk notat, hvor du først argumenterer for, at en afgørelse er ugyldig, og derefter skriver det stærkeste argument for, at den alligevel er gyldig.

\textbf{Kort svar:} Det stærkeste modargument er normalt ikke ``myndigheden gjorde alt rigtigt'', men en præcis indvending: der var hjemmel, oplysningen var uvæsentlig, parten kendte allerede forholdet, fejlen kunne ikke have påvirket resultatet, eller en særlig undtagelse gælder. At kunne formulere modargumentet gør hovedkonklusionen mere troværdig.

\section*{Samlet selvtest}

Hvis du kan besvare disse spørgsmål, har du forstået bogens grundmetode:
\begin{enumerate}
  \item Hvad er forskellen på en regel og en vurdering?
  \item Hvorfor kan en bekendtgørelse ikke uden videre skabe nye pligter?
  \item Hvad er forskellen på materiel ret og procesret?
  \item Hvorfor er et faktisk tab ikke automatisk et erstatningskrav?
  \item Hvad skal en proportionalitetsanalyse undersøge?
  \item Hvorfor skal AI-output kontrolleres mod primære kilder?
\end{enumerate}

